\documentclass[12pt,a4paper]{article}
%------------------------------------------------
% Кодировки и шрифты
%------------------------------------------------
\usepackage[utf8]{inputenc}
\usepackage[T2A]{fontenc}
\usepackage[russian]{babel}
\usepackage{amsmath,amssymb,mathtools}
%------------------------------------------------
% Оформление
%------------------------------------------------
\usepackage{geometry}
\geometry{left=20mm,right=20mm,top=20mm,bottom=20mm}
\usepackage{setspace}
\onehalfspacing
\raggedright
\setlength{\parindent}{0pt}
\setlength{\parskip}{0.7em}
\allowdisplaybreaks

\begin{document}

\noindent{\Large\textbf{Расстояние между точками на поверхности \(z=x^2+y^2\)}}
\vspace{0.4cm}

\textbf{Дано:}
\[
A(1,0,1), \qquad B(0,3,9).
\]
Поверхность:
\[
F(x,y,z)=x^2+y^2-z=0.
\]

\textbf{Найти:} длину геодезической, соединяющей точки \(A\) и \(B\).\\

%================================================
\textbf{\\1. Параметризация поверхности}
%================================================

Поверхность задана уравнением \(z=x^2+y^2\), введём полярные координаты:
\[
x=r\cos\varphi,\qquad y=r\sin\varphi\Rightarrow\ z=r^2\cos^2\varphi+r^2\sin^2\varphi.
\]
\(\cos^2\varphi+\sin^2\varphi=1\):
\[
z=r^2.
\]
\(\Rightarrow\)\[
\boxed{\vec r(r,\varphi)=(r\cos\varphi,r\sin\varphi,r^2)}
\]

Для точки \(A(1,0,1)\):
\[
r_A=\sqrt{1^2+0^2}=1,
\]
\[
\cos\varphi_A=1,\qquad \sin\varphi_A=0.
\]
\(\Rightarrow\)
\[
\boxed{A:(r,\varphi)=\left(1,0\right)}
\]

Для точки \(B(0,3,9)\):
\[
r_B=\sqrt{0^2+3^2}=3,
\]
\[
\cos\varphi_B=0,\qquad \sin\varphi_B=1.
\]
\(\Rightarrow\)
\[
\boxed{B:(r,\varphi)=\left(3,\frac{\pi}{2}\right)}
\]

\(\Rightarrow\) Гр. усл.:
\[
r(0)=1,\qquad r(S)=3,
\]
\[
\varphi(0)=0,\qquad \varphi(S)=\frac{\pi}{2}.
\]
\newpage
%================================================
\textbf{2. Вычисление первой квадратичной формы и метрики}
%================================================

Используем формулу первой квадратичной формы:
\[
ds^2=E\,dr^2+2F\,dr\,d\varphi+G\,d\varphi^2,
\]

\[
E=\langle\vec r_r,\vec r_r\rangle,
\qquad
F=\langle\vec r_r,\vec r_\varphi\rangle,
\qquad
G=\langle\vec r_\varphi,\vec r_\varphi\rangle.
\]

\[
\vec r_r=\frac{\partial\vec r}{\partial r}=(\cos\varphi,\sin\varphi,2r),
\]
\[
\vec r_\varphi=\frac{\partial\vec r}{\partial\varphi}=(-r\sin\varphi,r\cos\varphi,0).
\]

\[
E=\langle\vec r_r,\vec r_r\rangle.
\]
\[
E=\cos^2\varphi+\sin^2\varphi+4r^2.
\]
\(\Rightarrow\)
\[
\boxed{E=1+4r^2}
\]

\[
F=\langle\vec r_r,\vec r_\varphi\rangle.
\]
\[
F=-r\cos\varphi\sin\varphi+r\sin\varphi\cos\varphi=0.
\]
\(\Rightarrow\)
\[
\boxed{F=0}
\]

\[
G=\langle\vec r_\varphi,\vec r_\varphi\rangle.
\]
\[
G=r^2\sin^2\varphi+r^2\cos^2\varphi.
\]
\(\Rightarrow\)
\[
\boxed{G=r^2}
\]

\(\Rightarrow\) Первая квадратичная форма:
\[
\boxed{ds^2=(1+4r^2)dr^2+r^2d\varphi^2}
\]

Матрица метрики:
\[
\boxed{(g_{ij})=\begin{pmatrix} 1+4r^2 & 0 \\ 0 & r^2 \end{pmatrix}}
\]

Обратная матрица метрики:
\[
(g^{ij})=\begin{pmatrix} \dfrac{1}{1+4r^2} & 0 \\ 0 & \dfrac{1}{r^2} \end{pmatrix}.
\]\\

%================================================
\textbf{\\3. Функционал длины и уравнения Эйлера--Лагранжа}
%================================================

Используем §13.4, стр. 94 конспекта: уравнения Эйлера--Лагранжа для функционала
\[
S[\gamma]=\int_a^b L(u,\dot u,t)\,dt
\]
имеют вид
\[
\boxed{\frac{\partial L}{\partial u^i} - \frac{d}{dt} \left( \frac{\partial L}{\partial\dot u^i} \right)=0.}
\]

Согласно Теореме 23, §13.5, стр. 95--96, геодезические являются экстремалями функционала длины:
\[
ds^2=(1+4r^2)dr^2+r^2d\varphi^2 |: ds^2
\]

\[
1=(1+4r^2)\left(\frac{dr}{ds}\right)^2+r^2\left(\frac{d\varphi}{ds}\right)^2.
\]

\[
\dot r=\frac{dr}{ds},\qquad \dot\varphi=\frac{d\varphi}{ds}.
\]
\(\Rightarrow\) Функционал длины:
\[
S[\gamma]=\int \sqrt{(1+4r^2)\dot r^2+r^2\dot\varphi^2}\,ds.
\]
\(\Rightarrow\) Лагранжиан:
\[
L=\sqrt{(1+4r^2)\dot r^2+r^2\dot\varphi^2}.
\]
Поскольку \(s\) является натуральным параметром, скорость единичная: \(L=1\).
\(\Rightarrow\)\[\boxed{(1+4r^2)\dot r^2+r^2\dot\varphi^2=1.}\]
\newpage
%================================================
\textbf{4. Первый интеграл. Теорема Клеро}
%================================================

Используем Теорему 21 (Клеро), §13.2, стр. 91. Для поверхности вращения:
\[
\rho(s)\cos\alpha(s)=\mathrm{const},
\]
где \(\rho\) -- расстояние от точки до оси вращения, а \(\alpha\) -- угол между геодезической и параллелью.

В нашем случае \(\rho=r\). Единичный касательный вектор имеет компоненту в направлении параллели \(r\dot\varphi\).
\(\Rightarrow\)
\[
\cos\alpha=r\dot\varphi.
\]
Подставляем в теорему Клеро:
\[
r\cos\alpha=\mathrm{const},
\]
\[
r\cdot r\dot\varphi=\mathrm{const}.
\]
Обозначим постоянную через \(c\):
\[
\boxed{r^2\dot\varphi=c.}
\]

Получили первый интеграл геодезической. Его также можно получить непосредственно из уравнения Эйлера--Лагранжа. Действительно, лагранжиан не зависит явно от \(\varphi\):
\[
\frac{\partial L}{\partial\varphi}=0.
\]
Уравнение Эйлера--Лагранжа для \(\varphi\):
\[
\frac{d}{ds}\left(\frac{\partial L}{\partial\dot\varphi}\right)=0.
\]

\[
\frac{\partial L}{\partial\dot\varphi}=\frac{r^2\dot\varphi}{\sqrt{(1+4r^2)\dot r^2+r^2\dot\varphi^2}}.
\]
\(L=1\):
\[
\frac{d}{ds}(r^2\dot\varphi)=0.
\]

\[
r^2\dot\varphi=c.
\]
Совпадает с результатом Теоремы Клеро.
\(\Rightarrow\) Выражаем угловую скорость:
\[
\boxed{\dot\varphi=\frac{c}{r^2}.}
\]
\newpage
%================================================
\textbf{5. Выражение для радиальной скорости}
%================================================

Используем условие натуральной параметризации:
\[
(1+4r^2)\dot r^2+r^2\dot\varphi^2=1.
\]
Подставляем первый интеграл:
\[
(1+4r^2)\dot r^2+r^2\left(\frac{c}{r^2}\right)^2=1.
\]
\(\Rightarrow\)
\[
(1+4r^2)\dot r^2+\frac{c^2}{r^2}=1.
\]
\(\Rightarrow\)
\[
(1+4r^2)\dot r^2=1-\frac{c^2}{r^2}.
\]
\(\Rightarrow\)
\[
(1+4r^2)\dot r^2=\frac{r^2-c^2}{r^2}.
\]
\(\Rightarrow\)
\[
\dot r^2=\frac{r^2-c^2}{r^2(1+4r^2)}.
\]
\(\Rightarrow\)
\[
\dot r=\pm\frac{\sqrt{r^2-c^2}}{r\sqrt{1+4r^2}}.
\]
Рассматрим ветвь геодезической, на которой \(r\) возрастает от 1 до 3.\\

\\ \(\Rightarrow\) Выбираем положительный знак:
\[
\boxed{\frac{dr}{ds}=\frac{\sqrt{r^2-c^2}}{r\sqrt{1+4r^2}}.}
\]
\newpage
%================================================
\textbf{6. Дифференциальное уравнение геодезической}
%================================================

Найдём зависимость \(\varphi\) от \(r\):
\[
\frac{d\varphi}{dr}=\frac{d\varphi/ds}{dr/ds}.
\]
\(\Rightarrow\)
\[
\frac{d\varphi}{dr}=\frac{\dfrac{c}{r^2}}{\dfrac{\sqrt{r^2-c^2}}{r\sqrt{1+4r^2}}}.
\]
\(\Rightarrow\)
\[
\frac{d\varphi}{dr}=\frac{c}{r^2}\cdot\frac{r\sqrt{1+4r^2}}{\sqrt{r^2-c^2}}.
\]
\(\Rightarrow\)
\[
\boxed{\frac{d\varphi}{dr}=\frac{c\sqrt{1+4r^2}}{r\sqrt{r^2-c^2}}.}
\]
\\
%================================================
\textbf{7. Нахождение постоянной \(c\)}
%================================================
\[
d\varphi=\frac{c\sqrt{1+4r^2}}{r\sqrt{r^2-c^2}}\,dr.
\]
\(A:(1,0)\), \(B:\left(3,\frac{\pi}{2}\right)\).
\(\Rightarrow\)
\[
\int_{\varphi_A}^{\varphi_B}d\varphi=\int_{r_A}^{r_B}\frac{c\sqrt{1+4r^2}}{r\sqrt{r^2-c^2}}\,dr.
\]
\(\Rightarrow\)
\[
\int_0^{\pi/2}d\varphi=\int_1^3\frac{c\sqrt{1+4r^2}}{r\sqrt{r^2-c^2}}\,dr.
\]
\(\Rightarrow\)
\[
\boxed{\int_1^3\frac{c\sqrt{1+4r^2}}{r\sqrt{r^2-c^2}}\,dr=\frac{\pi}{2}.}
\]

Для взятия интеграла выполним замену переменной:
\[
u=\sqrt{r^2-c^2}.
\]
\(\Rightarrow\)
\[
u^2=r^2-c^2.
\]
\(\Rightarrow\)
\[
2u\,du=2r\,dr \implies r\,dr=u\,du.
\]

\[
1+4r^2 = 1+4(u^2+c^2) = 1+4c^2+4u^2.
\]
\(r\,dr=u\,du\):
\[
\frac{\sqrt{1+4r^2}}{r\sqrt{r^2-c^2}}\,dr = \frac{\sqrt{1+4c^2+4u^2}}{r \cdot u} \cdot \frac{u\,du}{r} = \frac{\sqrt{1+4c^2+4u^2}}{r^2}\,du.
\]
\(r^2 = u^2+c^2\):
\[
\frac{\sqrt{1+4r^2}}{r\sqrt{r^2-c^2}}\,dr = \frac{\sqrt{1+4c^2+4u^2}}{u^2+c^2}\,du.
\]
При \(r=1\):
\[
u_1=\sqrt{1^2-c^2}=\sqrt{1-c^2}.
\]
\(r=3\):
\[
u_2=\sqrt{3^2-c^2}=\sqrt{9-c^2}.
\]
\(\Rightarrow\)
\[
\int_{\sqrt{1-c^2}}^{\sqrt{9-c^2}}\frac{c\sqrt{1+4c^2+4u^2}}{u^2+c^2}\,du=\frac{\pi}{2}.
\]

Уравнение относительно неизвестной постоянной \(c\). Полученный интеграл является эллиптическим и не выражается через простые элементарные функции.\\

\(\Rightarrow\) Решаем уравнение численно. Вводим функцию:
\[
f(c)=\int_{\sqrt{1-c^2}}^{\sqrt{9-c^2}}\frac{c\sqrt{1+4c^2+4u^2}}{u^2+c^2}\,du-\frac{\pi}{2}.
\]
\(\Rightarrow\)\(f(c)=0\).
\(\Rightarrow\) Численное решение:
\[
c\approx0.62166256.
\]
\(\Rightarrow\) Проверка подстановкой:
\[
\int_1^3\frac{0.62166256\sqrt{1+4r^2}}{r\sqrt{r^2-0.62166256^2}}\,dr\approx1.57079633.
\]
При этом \(\frac{\pi}{2}\approx1.57079633\).
\(\Rightarrow\)
\[
\boxed{c\approx0.62166256.}
\]
\newpage
%================================================
\textbf{8. Вычисление длины геодезической}
%================================================

По определению натурального параметра:
\[
S=\int_A^B ds.
\]
Из выражения для радиальной скорости получаем:
\[
\frac{ds}{dr}=\frac{r\sqrt{1+4r^2}}{\sqrt{r^2-c^2}}.
\]
\(\Rightarrow\)
\[
S=\int_1^3\frac{r\sqrt{1+4r^2}}{\sqrt{r^2-c^2}}\,dr.
\]
\(\Rightarrow\) Подставляем найденное значение \(c\):
\[
S=\int_1^3\frac{r\sqrt{1+4r^2}}{\sqrt{r^2-0.62166256^2}}\,dr.
\]
\(\Rightarrow\) Сделаем замену:
\[
u=\sqrt{r^2-c^2}.
\]
\(\Rightarrow\)
\[
u^2=r^2-c^2.
\]
\(\Rightarrow\)
\[
2u\,du=2r\,dr.
\]
\(\Rightarrow\)
\[
r\,dr=u\,du.
\]
\(\Rightarrow\)
\[
r^2=u^2+c^2.
\]
\(\Rightarrow\)
\[
1+4r^2=1+4u^2+4c^2.
\]
\(\Rightarrow\)
\[
S=\int\frac{r\sqrt{1+4u^2+4c^2}}{u}\,dr.
\]
\(r\,dr=u\,du\):
\[
S=\int\sqrt{1+4u^2+4c^2}\,du.
\]
\(\Rightarrow\)
\[
a^2=1+4c^2.
\]
\(\Rightarrow\)
\[
S=\int\sqrt{4u^2+a^2}\,du.
\]
\(\Rightarrow\)
\[
\int\sqrt{4u^2+a^2}\,du=\frac{u}{2}\sqrt{4u^2+a^2}+\frac{a^2}{4}\ln\left(2u+\sqrt{4u^2+a^2}\right).
\]
\(\Rightarrow\)
\[
S=\left[\frac{u}{2}\sqrt{4u^2+a^2}+\frac{a^2}{4}\ln\left(2u+\sqrt{4u^2+a^2}\right)\right]_{u_1}^{u_2}.
\]

При \(r=1\):
\[
u_1=\sqrt{1-c^2}=\sqrt{1-0.62166256^2}\approx0.78328517.
\]
При \(r=3\):
\[
u_2=\sqrt{9-c^2}=\sqrt{9-0.62166256^2}\approx2.93488256.
\]
\(\Rightarrow\)
\[
a^2=1+4c^2=1+4\cdot0.62166256^2\approx2.54585734.
\]
\(\Rightarrow\)
\[
S=\left[\frac{u}{2}\sqrt{4u^2+2.54585734}+\frac{2.54585734}{4}\ln\left(2u+\sqrt{4u^2+2.54585734}\right)\right]_{0.78328517}^{2.93488256}.
\]
\(\Rightarrow\)
\[
\boxed{S\approx8.77926668.}
\]
\(\Rightarrow\) Длина найденной геодезической:
\[
\boxed{S_{\text{геод}}\approx8.7793.}
\]

%================================================
\textbf{9. Евклидово расстояние между точками}
%================================================

Для сравнения найдём обычное расстояние между точками в пространстве.\\
\(\Rightarrow\)
\[
\rho(A,B)=\sqrt{(x_B-x_A)^2+(y_B-y_A)^2+(z_B-z_A)^2}.
\]
\(\Rightarrow\)
\[
\rho(A,B)=\sqrt{(0-1)^2+(3-0)^2+(9-1)^2}.
\]
\(\Rightarrow\)
\[
\rho(A,B)=\sqrt{(-1)^2+3^2+8^2}.
\]
\(\Rightarrow\)
\[
\rho(A,B)=\sqrt{1+9+64}.
\]
\(\Rightarrow\)
\[
\boxed{\rho(A,B)=\sqrt{74}\approx8.6023.}
\]
\newpage
%================================================
\textbf{10. Итоговый ответ}
%================================================

\(\Rightarrow\)
\[
S_{\text{геод}}\approx8.7793.
\]
\(\Rightarrow\)
\[
\rho_{\text{Euclid}}=\sqrt{74}\approx8.6023.
\]
\(\Rightarrow\) Разница:
\[
\Delta S=8.77926668-8.60232527.
\]
\(\Rightarrow\)
\[
\boxed{\Delta S\approx0.17694.}
\]
\(\Rightarrow\)
\[
\boxed{\begin{aligned} S_{\text{геод}}&\approx8.7793,\\ \rho_{\text{Euclid}}&\approx8.6023. \end{aligned}}
\]

Геодезическая длина больше евклидова расстояния, поскольку геодезическая проходит по поверхности, а евклидово расстояние измеряется по прямому отрезку в пространстве.

\end{document}